\subsection{Repositorio del Proyecto (Monorepo)}

El código fuente integral de este proyecto se encuentra alojado de manera centralizada en un único repositorio de GitHub (monorepo). Esta estructura unificada alberga tanto el código fuente necesario para compilar esta documentación técnica en LaTeX (Game Design Document) como el proyecto completo y código fuente del videojuego desarrollado en el motor Unity.

\vspace{1em}

\noindent \textbf{Enlace al Repositorio Oficial:} \\
\url{https://github.com/Dennis290699/Project-Arrow}

\subsubsection*{Metodología de Trabajo y Control de Versiones}

El desarrollo del proyecto se rige bajo estrictas metodologías de control de versiones utilizando Git y GitHub como plataformas principales. Para garantizar la integridad del código, el correcto acoplamiento de nuevas características y la constante revisión por pares, el equipo de desarrollo emplea un flujo de trabajo riguroso basado en \textbf{Pull Requests (PRs)}.

El ciclo de desarrollo colaborativo se estructura de la siguiente manera:

\begin{itemize}
    \item \textbf{Gestión de Ramas (Branching):} Cada integrante del equipo crea ramas específicas de características (\textit{feature branches}) derivadas de la rama principal para aislar el desarrollo de nuevas mecánicas, redacción de documentación o corrección de errores.
    \item \textbf{Pull Requests (Solicitudes de Extracción):} Una vez culminada la tarea asignada, el autor emite un \textit{Pull Request} hacia la rama principal. Esto actúa como un punto de control formal y provee trazabilidad histórica sobre cada adición al código fuente.
    \item \textbf{Revisión por Pares (Peer Review) y Aprobación:} Ningún cambio es integrado de manera directa. Los \textit{Pull Requests} son sometidos a la escrutinio del resto de los integrantes, quienes pueden realizar retroalimentación, solicitar modificaciones o brindar su aprobación final (\textit{Approve}).
    \item \textbf{Fusión (Merging):} Únicamente tras superar el proceso de revisión, debate y validación por parte del equipo, los cambios son fusionados (\textit{merged}) a la rama principal, asegurando un entorno estable y minimizando los conflictos de integración.
\end{itemize}

Este flujo de trabajo profesional facilita la colaboración asíncrona y simultánea entre los integrantes, y garantiza un historial de cambios completamente auditable a lo largo de todo el ciclo de vida de la creación del videojuego y su documentación.
